\clearpage\thispagestyle{empty}
\pdfbookmark[0]{Quellen zu Gauss an Schumacher (17. Januar 1844)}{cometletterjan1844-sources}
\begin{center}\textbf{Quellen zu Gauss an Schumacher (17. Januar 1844)}\end{center}
\noindent Diese KI-Übertragung ins heutige Deutsch folgt dem Text auf S.~464--465 der \href{https://gdz.sub.uni-goettingen.de/id/PPN236007467}{\textit{Werke}, Band~VI} (Göttingen, 1874). Carl Friedrich Gauß schrieb den Brief am 17. Januar 1844 in Göttingen an Heinrich Christian Schumacher. Das Heft des Erstdrucks ist auf den 8. Februar 1844 datiert. \textit{Astronomische Nachrichten}, Band~XXI, Nr.~495, \href{https://api.digitale-sammlungen.de/iiif/image/v2/bsb10049446_00126/full/1800,/0/default.jpg}{Sp.~235--236}.
\par\smallskip\noindent Die Quellenangabe der Werkausgabe lautet \textit{Astronomisehe}. Ihre Tabelle nennt den 13. Dezember 1842, die des Erstdrucks den 13. Dezember 1843. Die Werkausgabe spricht von mehreren Chronometern, der Erstdruck von einem. Die Übertragung behält die Lesarten der Werkausgabe und die beiden Punkte bei den Sekundenangaben bei. Der Erstdruck enthält in \href{https://api.digitale-sammlungen.de/iiif/image/v2/bsb10049446_00127/full/1800,/0/default.jpg}{Sp.~237--238} auch die Ephemeride und die zweiten Elemente. Sie gehören nicht zum Umfang des gesammelten Auszugs. KI-Übertragung und KI-Satz.
\clearpage
